% 本文件由 LaTeX 排版 Agent 根据有效 translation.json 自动生成。
% author=Augustine of Hippo; work_id=1404; title=行为录或驳福图纳图斯辩

\chapter{\WorkTitle{行为录或驳福图纳图斯辩}}

% structure: p0001 source=h1 target=omitted reason=page-title-already-rendered-by-parent

% structure: p0002 source=h2 target=section reason=relative-heading-rank; base=h2

\section{第一日的辩论}

九月五日，最尊贵之人\Person{阿卡狄乌斯·奥古斯都}（第二次执政）与\Person{鲁菲努斯}任执政官之际，在\Place{希波勒吉乌斯}城的索西乌斯浴场，当着民众之面，举行了一场反对摩尼教长老\Person{福图纳图斯}的辩论。\par

1. 我现在视之为错误的，原是我先前视之为真理的。我渴望听你们在场的人说，我的猜想是否正确。首先，我认为相信全能天主在任何部分是可侵犯的、可玷污的、可朽坏的，是最大的错误。我知道这是你们的异端所肯定的，虽然并非我现在所用的措辞；因为当你们被问及时，你们承认天主是不朽坏的、完全不可侵犯的、不可玷污的；但当你们开始阐述你们体系的其余部分时，我们就不得不宣称祂是可朽坏的、可穿透的、可玷污的。因为你们说，另一个黑暗种族，无论是什么，反抗了天主的国度；但全能天主看见祂的领域受到毁灭和荒凉的威胁，除非祂对敌对种族作出某种抵抗并抵挡它，就派遣了这个德能，从它与邪恶和黑暗种族的混合中，世界得以形成。因此，善灵魂在这里劳作、服役、犯错、被腐蚀：好使他们看到需要一位解救者，祂应净化他们脱离错误，将他们从与邪恶的混合中释放，并解救他们脱离奴役。我认为相信全能天主曾害怕任何敌对种族，或不得不将我们投入苦难，是不敬的。\par

\Person{福图纳图斯}说：因为我知道你曾在我们中间，即作为附从者生活在摩尼教徒中，这些是我们信仰的原则。现在需要考虑的是我们的生活方式，那些被诬告的罪行，我们因此受到虐待。因此，请在场的善人们从你那里听到，我们所被指控的、被别人扔在我们脸上的这些事是真还是假。因为通过你的教导，你的阐述和解释，他们若经你阐述，就能对我们的生活方式获得更正确的认识。\par

2. \Person{奥古斯丁}说：我曾在你中间，但信仰与道德是不同的问题。我提议讨论信仰。但如果在场的人更愿意听关于道德的问题，我也不拒绝那个问题。\par

\Person{福图纳图斯}说：我首先希望，藉着一位胜任者（甚至现在对我也是胜任者）的见证，并在基督公义审判者将来审查之前，在你那被我们玷污的良心上为自己辩护：他是否看到在我们中间，或者他自己通过模仿实行了那些现在被扔到我们脸上的事？\par

3. \Person{奥古斯丁}说：你召唤我去做别的事，我原本提议讨论信仰。但关于你们的道德，只有你们的\OriginalTerm{选民}{Elect}才能完全知道。但你知道我不是你们的\OriginalTerm{选民}{Elect}，而是\OriginalTerm{听者}{Auditor}。因此，虽然我参加了你们的祈祷聚会，如你所问（你们彼此之间是否有单独的祈祷聚会，只有天主和你们自己知道），但在我在场的祈祷聚会中，我没有看到任何可耻的事发生；只是后来我所学习并赞同的信仰遭到指责，并且你们面向太阳举行仪式。除此之外，我在你们的聚会上没有发现任何新东西；但任何人提出关于道德的问题，都是针对你们的\OriginalTerm{选民}{Elect}。但你们这些\OriginalTerm{选民}{Elect}在你们中间做什么，我没有办法知道。因为我常听你们说你们领受\OriginalTerm{圣体}{Eucharist}。但领受的时间向我隐瞒，我怎能知道你们领受什么？所以，如果你愿意，请将道德的问题留给你们的\OriginalTerm{选民}{Elect}讨论，如果能够讨论的话。你给我的信仰是我今天不赞同的。我提议讨论这个。请对我的提议作出回应。\par

\Person{福图纳图斯}说：而我们的宣认正是如此：天主是不朽坏的、光明的、不可接近的、不可把握的、不受苦难的，祂居住在自己永恒的光明中，从祂那里没有可朽坏的东西发出，黑暗、魔鬼、撒旦或任何敌对之物都不能在祂的国度中找到。但祂派遣了一位与自己相似的救主；从世界奠基时诞生的圣言，在形成世界之后，来到人间；祂按照自己神圣的旨意拣选了配得上自己的灵魂，以属天的命令圣化，以属天事物的信仰和理性浸润；那些灵魂将在祂的领导下，按照祂的神圣应许返回天主的国度，祂说：\BibleQuote{我是道路、真理、门；}并且\BibleQuote{除非经过我，没有人能到父那里去。}因为这些事我们相信，因为否则，即通过另一位中保，灵魂不能返回天主的国度，除非他们找到祂作为道路、真理和门。因为祂亲口说：\BibleQuote{看见我的，也看见了父；}\BibleRef{若 14:8} 并且\BibleQuote{凡信我的人永远不死，却已出死入生，不会进入审判。}\BibleRef{若 5:24} 我们相信这些事，这就是我们信仰的理由，并且按照我们心智的力量，我们努力遵行祂的诫命，追随这位天主圣三、圣父、圣子、圣神的一信。\par

4. \Person{奥古斯丁}说：那么，那些你承认藉着基督从死亡进入生命的灵魂，是什么原因导致它们被投入死亡？\par

\Person{福图纳图斯}说：那现在请继续并予以反驳，如果除了天主之外没有别的。\par

5. \Person{奥古斯丁}说：不，请你回答所问的问题：是什么原因使这些灵魂死亡？\par

\Person{福图纳图斯}说：不，但请你回答，除了天主之外是否有别的东西，或者万有都在天主内。\par

6. \Person{奥古斯丁}说：我可以这样回答：主愿意我知道，天主不能遭受任何必然性，也不能在任何部分被侵犯或被朽坏。这一点，既然你也承认，我问：祂出于什么必然性派遣了这些你所说的藉着基督返回的灵魂？\par

福图纳图斯说：你所说，即天主迄今为止向你启示祂是不可朽坏的，正如祂也向我启示的一样，那么必须探究灵魂如何以及为何进入这个世界，以致现在天主应当通过祂唯一且与祂相似的圣子来解救他们脱离这个世界，若除祂以外别无他物？\par

7. 奥古斯丁说：我们不应使在场的诸位，即那些显要人士失望，也不应从所提议讨论的问题转向另一个。因此我们双方都承认，都赞同：天主是不可朽坏、不可侵犯的，绝不可能受苦。由此推出，你们的异端是虚假的，它说：当天主看到荒凉和毁灭威胁祂的国度时，祂派遣一种力量去与黑暗种族作战，而由此混杂中我们的灵魂正在受苦。我的论证是简短的，并且我认为对任何人都是完全清楚的。如果天主不可能因黑暗种族而受苦，因为祂是不可侵犯的，那么祂无故派遣我们到这里来受苦。但如果有什么能够受苦，那它就不是不可侵犯的，而你们欺骗了那些听你们说天主是不可侵犯的人。因为你们的这个异端在你们阐释其余部分时否认了这一点。\par

福图纳图斯说：我们持有宗徒保禄教导我们的心意，他说：\BibleQuote{你们该怀有基督耶稣所怀有的心情：祂虽具有天主的形体，却没有以自己与天主同等为应当把持不舍的，反而空虚自己，取了奴仆的形体，成了人的模样；以人的形态出现后，祂贬抑自己，听命至死。}\BibleRef{斐 2:5}因此，我们关于自己有这同样的心意，正如我们关于基督所怀有的：祂虽具有天主的形体，却听命至死，是为了显示我们灵魂的相似。正如祂在自己身上显示了死亡的相似，并藉着从死者中被举起而显示祂来自父，同样我们想这对于我们的灵魂也是如此，因为藉着祂我们终将能够脱离这死亡——这死亡要么是与天主无关的，要么如果它属于天主，祂的仁慈就停止了，解放者的名号以及祂解放的作为也停止了。\par

8. 奥古斯丁说：我问我们是如何进入死亡的，而你说我们如何从死亡中被解放。\par

福图纳图斯说：所以宗徒说，我们应当怀有基督向我们显示的那同样的心意。如果基督曾处于苦难和死亡之中，我们也是如此。\par

9. 奥古斯丁说：众所周知，公教信仰是这样的：我们的主，即天主的德能和智慧\BibleRef{格前 1:24}，以及万物藉之而造、没有祂便无物受造的圣言\BibleRef{若 1:3}，取了人来解救我们。在祂所取的人性中，祂展现了你们所提及的那些事。但现在我们问的是天主本身的实体和不可言喻的尊威，是否有什么能伤害它或者不能。因为如果有东西能伤害它，祂就不是不可侵犯的。如果没有什么能伤害天主的实体，那么黑暗种族对它有什么可做的呢？你们说在创世之前天主就与它作战；在那个战争中，你们声称我们——即现在明显需要解救者的灵魂——已经与一切恶混杂并陷入死亡。因为我回到那个非常简短的陈述：如果祂能被伤害，祂就不是不可侵犯的；如果不能，那么祂派遣我们来到这里受这些苦就是残忍的。\par

福图纳图斯说：灵魂属于天主，还是不属于？\par

10. 奥古斯丁说：如果你不回答我的问题，而由我来被询问，这是否公平？如果是公平的，我将回答。\par

福图纳图斯说：灵魂是否独立行动？我问你这个问题。\par

11. 奥古斯丁说：我确实会告诉你所问的；只是请记住这一点：当你拒绝回答我的问题时，我却回答了你的。如果你问灵魂是否来自天主，这确实是一个大问题；但无论灵魂是否来自天主，我对灵魂的回答是：它不是天主；天主是一回事，灵魂是另一回事。天主是不可侵犯、不可朽坏、不可渗透、不可玷污的，祂在任何部分都不可能被败坏，也没有任何部分能受伤害。但我们也看到灵魂是有罪的，处于痛苦中，寻求真理，需要解救者。灵魂的这种变化状态向我表明灵魂不是天主。因为如果灵魂是天主的实体，那么天主的实体犯错，天主的实体被败坏，天主的实体被侵犯，天主的实体被欺骗——这是亵渎的说法。\par

福图纳图斯说：因此你否认灵魂来自天主，只要它服事罪、恶习和属世之事，并被错误引导，因为不可能天主或祂的实体遭受这事。因为天主是不可朽坏的，祂的实体是无玷污和神圣的。但这里问的是你：灵魂是否来自天主？我们承认这一点，并从救主的来临、祂神圣的宣讲、祂的拣选中显示出来；祂怜悯灵魂，据说灵魂按祂的旨意而来，为的是祂能解救它脱离死亡，引领它到永久的荣耀，并把它归还给父。但你对灵魂持什么观点和盼望？它是来自天主吗？你否认灵魂来自天主，那么天主的实体能否不受任何情欲的影响？\par

12. 奥古斯丁说：我否认灵魂是天主的实体，意思是灵魂不是天主；但我仍坚持灵魂来自作为其创作者的天主，因为它是由天主造的。造物主是一回事，受造物是另一回事。那造物的绝不能有丝毫可朽坏，而祂所造的绝不能与造祂的相等。\par

福图纳图斯说：我也没有说过灵魂像天主。但是因为你说了灵魂是造出来的东西，并且除天主之外别无他物，我问：那么天主从哪里发明了灵魂的实体？\par

13. 奥古斯丁说：请记住，我回应你的质询，但你并未回应我的。我说灵魂是由天主所造，正如其他一切由天主所造的事物一样；而在全能天主所造之物中，灵魂被赋予首要地位。但若你问天主从何造灵魂，请记住你我都一致承认天主是全能的。但若他需要借助任何质料才能造他所愿之物，他就不再是全能的。由此推论，按照我们的信德，天主藉其圣言和智慧所造的一切，都是从无中造的。因为经上记载：\Quote{祂下令，他们就造成；祂命令，他们就受造。}\par

福尔图纳图斯说：万物皆因天主的命令而存在吗？\par

14. 奥古斯丁说：我如此相信，但仅限于一切受造之物。\par

福尔图纳图斯说：作为受造之物，它们是一致的；但由于它们彼此不相适合，因此由此得出，并非只有一个实体，尽管它们出于同一个\Quote{一}的秩序而来到这世界的组合和构造中。但在事物本身中明显可见，黑暗与光明、真理与虚假、死亡与生命、灵魂与身体，以及其他在名称和外观上彼此不同的事物之间，并无相似之处。我们的主有理由说：\Quote{那棵我天父未栽种的树，必被连根拔起扔进火里，因为它不结好果子。}并且那树已被连根拔起。因此，从事物的理性出发，这世界确实有两种实体，它们在形式和名称上是一致的，其中一种属于有形体的本性，而另一种则是全能父的永恒实体，我们相信这实体就是天主的实体。\par

15. 奥古斯丁说：那些使你困扰、使我们思想对立的相反之事，是因我们的罪、即人的罪而发生的。因为天主造了万物，都是好的，并妥善安排了它们；但他并没有造罪，唯独我们自愿的罪才被称为恶。另有另一种恶，即罪的惩罚。既然有两种恶——罪和罪的惩罚，罪与天主无关；罪的惩罚属于报复者。因为正如天主是善的，创造了万物，他亦是公义的，惩罚罪。既然万物都以最好的方式被安排——那些如今在我们看来相反的事，恰是堕落之人不肯遵守天主的律法之后果。因为天主赋予人内理性灵魂以自由意志。因为这样我们才能有功绩，若我们是自愿为善而非出于必然。因此，既然我们应该出于自愿而非出于必然成为善，天主就有必要赋予灵魂自由意志。对于这服从他律法的灵魂，他使万物毫无逆境地臣服于它，以便天主所造的其他事物都服务于它，如果灵魂本身也愿意服事天主的话。但如果它拒绝服事天主，那些曾服事它的事物就会变成对它的惩罚。因此，若万物都由天主妥善安排且是善的，天主也不受恶之苦。\par

福尔图纳图斯说：他不受苦，而是阻止恶。\par

16. 奥古斯丁说：那么恶是从谁那里使他受苦的呢？\par

福尔图纳图斯说：这正是我的论点：他希望阻止恶，不是鲁莽地，而是藉着能力和预知。但否认恶是独立于天主的，因为可以指出有其他诫命是违反他旨意而行的。除非有相反的情况，否则不会引入诫命。自由生活的能力只有存在堕落的情况下才被赐予，正如宗徒的论述所说：\Quote{你们从前因过犯和罪恶是死的，那时你们行事为人，随从今世的风俗，顺服空中掌权者的首领，就是现今在悖逆之子心中运行的邪灵。我们从前也都在他们中间，放纵肉体的私欲，随着肉体和心中所喜好的去行，本为可怒之子，和别人一样。然而天主既有丰富的怜悯，因他爱我们的大爱，当我们死在过犯中的时候，便叫我们与基督一同活过来（你们得救是本乎恩宠）。他又叫我们与基督耶稣一同复活，一同坐在天上，要将他极丰富的恩宠，就是他在基督耶稣里向我们所施的恩慈，显明给后来的世代看。你们得救是本乎恩宠，也因着信德，这并不是出于自己，乃是天主的恩赐；也不是出于行为，免得有人自夸。我们原是他的工作，在基督耶稣里造成的，为要叫我们行善，就是天主所预备叫我们行的。所以你们应当记念，你们从前按肉体是外邦人，是称为没受割礼的，这名原是那些凭人手在肉身上称为受割礼之人所起的。那时你们与基督无关，在以色列国民以外，在所应许的诸约上是局外人，并且活在世上没有指望，没有天主。你们从前远离天主的人，如今却在基督耶稣里，靠着他的血，已经得亲近了。因他使我们和睦，将两下合而为一，拆毁了中间隔断的墙，而且以自己的身体废掉冤仇，就是那记在律法上的规条，为要将两下藉着自己造成一个新人，如此便成就了和睦。既在十字架上灭了冤仇，便藉这十字架使两下归为一体，与天主和好了。并且来传和平的福音给你们远处的人，也给那近处的人。因为我们两下藉着他被一个圣神所感，得以进到父面前。}\par

17. 奥古斯丁说：你所引用的宗徒的这段话，如果我没有弄错的话，非常有力地支持了我的信仰，反对了你的信仰。首先，因为自由意志本身（我曾说过，灵魂犯罪的可能性依赖于它）在这里得到了充分的表达，当提到罪时，并说我们藉着耶稣基督与天主和好。因为犯罪使我们与天主对立；但遵行基督的诫命使我们与天主和好；这样，我们这些死于罪恶的人，因遵守祂的诫命而得以存活，并在同一圣神内与祂和好，而从前我们因不遵守祂的诫命而与祂隔绝；正如我们的信仰中关于最初受造的人所教导的。因此，我请问你，根据刚才所读的经文，如果\OriginalTerm{相反本性}{contrary nature}强迫我们去做我们所做的事，我们怎么可能有罪？因为凡是被本性强迫去做任何事的人，并不犯罪。但犯罪的人，是凭自由意志犯罪。既然如此，如果我们没有做任何恶事，只有\OriginalTerm{黑暗种族}{race of darkness}做了，悔改又为什么命令我们去做？同样，我问你，罪的赦免是赐给我们，还是赐给黑暗种族？如果是赐给黑暗种族，那么他们的种族也将与基督一同为王，接受罪的赦免；但如果是赐给我们，那就表明我们自愿犯了罪。因为赦免一个没有作恶的人是完全荒谬的。但一个人若没有照自己的意志行事，他就没有作恶。因此，今天那灵魂应许自己罪的赦免和与天主的和好，如果它停止犯罪，并悔改过去的罪；如果它按照你们的信仰回答，并说：\Quote{我在什么事上犯了罪？我在什么事上有罪？你为什么将我逐出你的疆域，让我去与某个种族作战？我被人践踏，我被混合，我被腐蚀，我精疲力竭，我的自由意志没有被保存。你知道我得以保存的必然性：你为什么将我所受的创伤归咎于我？你为什么强迫我悔改，而你是造成我创伤的原因？你知道我遭受了什么，黑暗种族对我做了什么，你是那不能受害的作者，却想拯救那些不受损害的疆域，你把我推入了这些苦难。如果我真是一部分来自你，从你的怀中发出，如果我是来自你的国度、来自你的口，我就不该在这个黑暗种族中遭受任何事，这样，我若是不受腐蚀的，那个种族就应被制服，如果我是主的一部分。但现在既然除非通过我的腐蚀，它就不能被控制，我怎么能被称为你的一部分，或者你保持不可侵犯，或者你不残忍，因为我希望我为那些疆域受苦，而那些疆域绝不可能被那个黑暗种族所伤害？}请回答这个问题，并恻隐地向我解释，宗徒怎么说\BibleQuote{我们生来都是义怒之子}，他说，这些人与天主和好。因此，如果他们生来都是义怒之子，你怎么说灵魂生来是天主的女儿和一部分？\par

福尔图纳图斯说：如果宗徒关于灵魂说了\BibleQuote{我们生来都是义怒之子}，那么灵魂就被宗徒的口与天主疏远了。从这个论证中，你只是表明灵魂不属于天主，因为宗徒说\BibleQuote{我们生来都是义怒之子}。但是，如果考虑到宗徒\BibleRef{罗 11:1}被律法所束缚（正如他自己作证，出自亚巴郎的后裔），那么他就是在肉身方面说的，即我们（\Emph{i.e.}犹太人）生来都是义怒之子，如同其余的人一样。但他表明灵魂的实体出于天主，并且灵魂除了通过那位导师，即基督耶稣，不能与天主和好。因为仇恨被杀死后，灵魂似乎在天主眼中不配存在。但是，灵魂被派遣，我们承认，是由全能的天主，既从祂而出，又被派遣来封印祂的旨意。同样，我们也相信基督救主从天降下，为承行父的旨意。父的旨意就是，这种仇恨被杀死后，释放我们的灵魂脱离同样的仇恨；这种仇恨若不曾与天主对立，在统一中就不可称为仇恨，也不能说被杀死，或在有生命的地方发生。\par

18. 奥古斯丁说：请记住，宗徒说我们因我们的生活方式而与天主疏远。\par

福尔图纳图斯说：我坚持，有两种实体。在\OriginalTerm{光明实体}{substance of light}中，如上所述，天主应被视为不可朽坏的；但有一种\OriginalTerm{黑暗相反本性}{contrary nature of darkness}，我今天也承认，它被天主的德能所征服，并且基督已被派遣作为救主来恢复我，正如先前同一宗徒所说的。\par

19. 奥古斯丁说：我们应该在理性基础上讨论对两种本性的信仰，这是由那些听我们讲话的人所要求的。但既然你又转向了圣经，我就下降到它们之中，并要求任何事都不要被忽略，以免使用某些陈述给那些不熟悉圣经的人带来困惑。因此，让我们思考宗徒在《罗马书》中的一句话。因为在第一页上就有强烈反对你的言论。他说：\BibleQuote{保禄，耶稣基督的仆人，蒙召作宗徒，被选拔为传天主的福音，这福音是天主先前藉着祂的先知在圣经中所预许的，关于祂的儿子，按肉身说，是从达味的后裔生的，按圣洁的灵说，由于从死者中复活，以大能显明是天主的儿子，我们的主耶稣基督。}\BibleRef{罗 1:1}我们看到，宗徒教导我们关于主耶稣基督，在肉身之前，他已由天主的德能所预定，按肉身说，是由达味的后裔生的。既然你们一直否认并将永远否认这一点，你们为什么如此迫切地要求根据圣经来讨论呢？\par

\Person{福图纳图斯}说：\Quote{你断言基督按肉身是\Person{达味}的后裔，然而应当断言祂是童贞女所生，\BibleRef{依 7:14}并应被尊为天主子。因为除非如同出自圣神的可被视为神，如此出自肉身的也可被知为肉，这是不可能的。\BibleRef{若 3:6}与福音的权威相反，其中说：\BibleQuote{血肉不能承受天主的国，朽坏也不能承受不朽坏。}\BibleRef{格前 15:50}}\par

此时听众发出喧哗，他们希望辩论在理性基础上进行，因为他们看到\Person{福图纳图斯}不愿意接受\OriginalTerm{宗徒书信集}{Codex}中记载的一切。然后大家开始到处进行零星讨论，直到\Person{福图纳图斯}说天主的圣言被束缚在黑暗种族中。对此，在场的人表达了恐惧，会议结束。\par

% structure: p0043 source=h2 target=section reason=relative-heading-rank; base=h2

\section{第二天的辩论}

次日，再次召来公证人，辩论如下：\par

\Person{福图纳图斯}说：\Quote{我说全能天主从自身不产生任何恶，属于祂的事物保持不朽，因为它们起源于并生于一个不可侵犯的源头；但其他相反的事物，它们存在于这个世界，并非从天主流出，也不是以天主为作者出现在这个世界；也就是说，它们并非源自天主。因此我们相信恶的事物与天主无关。}\par

\Person{奥斯定}说：\Quote{我们的信仰也是如此：天主不是恶的源头，也没有造任何恶的本性。但既然我们双方都同意天主是不朽坏、不可玷污的，那么谨慎而有信德的人应该考虑，哪一种信仰更纯洁、更配得上天主的尊威：是那种主张天主的权能、或天主的某一部分、或天主的圣言可以被改变、侵犯、腐化、束缚的信仰；还是那种主张全能天主及其全部本性和实体在任何部分都永不可被腐化，而恶的存在是由于灵魂自愿的罪恶，天主赋予灵魂自由意志。如果天主没有赋予自由意志，就不会有公义的惩罚审判，不会有义行的功德，不会有悔改罪过的神圣训导，也不会有天主通过我们的主耶稣基督赐予我们的罪过的赦免本身。因为不是自愿犯罪者，根本不犯罪。我想这对所有人都是明白清楚的。因此，如果我们按照自己的功过在天主所造的事物中遭受一些不便，我们不应为此烦恼。因为正如祂是善的，所以祂创造一切；同样祂是公义的，所以祂不宽恕罪过。这些罪过，如我所说，除非我们有自由意志，否则就不会是罪。因为如果有人，比如说，被某人捆住其他肢体，而用手违背其本意写下虚假的东西，我问，如果这在法官面前公开，他能否判这个人犯有虚假罪？因此，如果显而易见，没有自由意志的行使就没有罪，那么我想听听，你所称之为天主的某一部分、或权能、或圣言、或别的什么的灵魂，做了什么恶，以致被天主惩罚，或悔改罪过，或得到赦免的功德，既然它根本没有犯罪？}\par

\Person{福图纳图斯}说：\Quote{我关于实体提出，天主只应被视为善的事物的创造者，但作为恶的事物的惩罚者，原因是恶的事物并非来自祂。因此我理由充分认为这一点，并且天主惩罚恶的事物是因为它们不是来自祂。但如果它们来自祂，要么祂会给予它们犯罪的许可，如你所说天主赋予自由意志，那么祂就会被发现是我过错的参与者，因为祂是我过错的作者；要么祂不知道我将成为什么，祂留下我这祂并没有使之配得上祂自己的人。因此这就是我所提出的，现在我问的是，天主是否创造了恶？以及祂是否亲自设立了恶的终结。因为从这些事可以看出，而且福音信仰教导说，我们所说的那些由天主自己作为创造主天主所造的事物，作为由祂创造和生出的，应被视为不朽坏的。这些也是我提出的属于我们信仰的，并且可以被你在我们那个认信中所确认，而不损害基督教信仰的权威。因为我无法表明我正确地信仰，除非我通过圣经的权威来确认那个信仰，因此这就是我所暗示、所说的。要么如果恶的事物与天主一起作为它们的作者出现在世界，请你自己这样说；要么如果相信恶的事物不是来自天主的观点是正确的，那么在场者的考量也应当尊重并接受它。我谈论的是实体，而非居住在我们内的罪。因为如果我们认为制造过错的东西没有起源，我们就不会被强迫去犯罪或犯错。因为我们不情愿地犯罪，并被一个与我们自身相反且敌对的实体所强迫，因此我们遵循对事物的知识。灵魂通过这种知识得到劝诫并恢复至原始记忆，认识到它自身从何起源，它居住在何种恶中，通过何种善功再次纠正它不情愿犯罪之处，它可以通过对其过错的纠正，为了善功的缘故，为自己确保与天主和好的功德，我们的救主是它的作者，祂也教导我们行善并远离恶。因为你要求我们相信，不是通过某种相反的本性，而是通过人自己的选择，人要么事奉正义，要么陷入罪恶；既然没有相反的种族存在，如果灵魂（如你所说天主赋予其自由意志），被安置在身体中，单独居住，它就会无罪，也不会陷入罪恶。}\par

21. \Person{奥古斯丁}说：我说，若非出于自己的意志，就不是罪；因此也有赏报，因为我们出于自己的意志行善。或者，若那出于不情愿而犯罪的人应受惩罚，那么出于不情愿而行善的人也应当得到赏报。但谁怀疑赏报只赐给那些出于善意而行某事的人呢？由此我们知道，惩罚也是加给那些出于恶意而行某事的人。但既然你把我拉回到原初的本性与实体，我的信德是这样的：全能的天主——这尤其需要注意并牢记在心——全能的天主创造了美善的事物。但祂所创造的事物不可能像创造它们的那一位那样。因为相信工程与工匠平等、受造物与造物者平等，是不义且愚蠢的。因此，若怀着虔敬相信天主创造了所有美善的事物（然而祂远超过它们，远优越于它们），那么邪恶的起源与首脑就是罪，正如宗徒所说：\BibleQuote{贪欲乃是一切邪恶的根源；有些人随从贪欲，就离弃了信德，使自己深受许多痛苦。}(\BibleRef{弟前 6:10}) 因为若你寻求一切邪恶的根源，你有宗徒说贪欲是一切邪恶的根源。但我不能寻求根源的根源。或者，若有另一种邪恶，其根源不是贪欲，那么贪欲就不是一切邪恶的根源。但若贪欲真是一切邪恶的根源，那么我们再寻求其他种类的邪恶就是徒劳的。至于你所引进的那个相反本性，既然我已经回答了你的反对意见，我请求你俯允告诉我：它是否完全邪恶？是否除了它之外不可能有罪？是否唯独它应受惩罚，而非那没有犯过罪的灵魂？但若你说唯独这相反本性应受惩罚，而非灵魂，那么我问，对于那被命令的、所赐予的补赎，是给谁的呢？若灵魂被命令补赎，罪就是出自灵魂，并且灵魂是自愿犯罪的。因为若灵魂被迫作恶，那么它所行的就不是恶。说黑暗种族犯了罪，而我为罪过补赎，这难道不是愚蠢且最荒谬的吗？说黑暗种族犯了罪，而罪的赦免赐予了我——我按照你们的信德完全可以这样说：我做了什么？我犯了什么？我曾与祢同在，我处在完好状态，未受任何玷污。祢派遣我到这里，祢忍受了必然性，当巨大的玷污与荒凉威胁着祢的领域时，祢保护了它们。既然祢知道那压迫我的必然性，使我无法呼吸，我无法抵抗；为何祢指责我好像我在犯罪？或为何祢应许罪的赦免？请毫不回避地答复这一点，若你愿意，正如我已答复了你。\par

\Person{福尔图纳图斯}说：我们说，灵魂被相反本性强迫而干犯过犯，对于这过犯，你认为除了那住在我们内的邪恶之外没有别的根源；因为毫无疑问，除了我们的身体之外，邪恶之物遍布全世界。因为不仅那些我们身体内之物遍布全世界，并以善的名称为人所知；一个邪恶的根源也内附着。因为你的尊贵说，这住在我们身体内的贪欲是邪恶的根源；既然在我们身体之外没有对邪恶的贪欲，那么相反本性就从那里遍及全世界。因为宗徒把那个，即贪欲，指明为邪恶的根源，而不是你称为一切邪恶之根源的那一个邪恶。但贪欲——你曾说它是一切邪恶的根源——不是以一种方式被理解，仿佛只限于住在我们身体内的那个；因为毫无疑问，这住在我们内的邪恶是从一个邪恶的作者降下来的，并且这个你称为根源的，只是邪恶的一小部分，所以它不是根源本身，而是那无处不在的邪恶的一小部分。我们的主将这个根源和树称为邪恶的树，从不结好果子，是祂的父没有栽种的，理应被连根拔起丢在火里。因为如你所说，罪应归咎于相反本性，那本性属于邪恶；而灵魂的罪是这样：若在救主的警告和祂有益的教训之后，灵魂将自己与其相反的和敌对的种族分隔开来，并以更纯洁的事物装饰自己；否则它不能恢复到自己原来的实体。因为经上说：\BibleQuote{我若没有来对他们说话，他们就没有罪；但现在我来了，对他们说了话，他们却拒绝相信我，他们的罪就无可推诿了。}(\BibleRef{若 15:22}) 由此非常清楚，补赎是在救主来临之后，并在这事物的认识之后赐下的，灵魂借此可以如同在神圣的泉源中洗净来自全世界的以及那灵魂所住的身体的污秽和恶习，从而恢复它曾离开的天国。因为宗徒说：\BibleQuote{肉性的思念是与天主为敌的；不服天主的法律，也是不能服。}(\BibleRef{罗 8:7}) 因此从这些事明显看出，善良的灵魂看似不是自愿犯罪，而是由于那不服天主法律的行为。同样也接着说：\BibleQuote{肉情和圣神相争，圣神和肉情相争，使你们不能作你们所愿意的。}(\BibleRef{迦 5:17}) 又说：\BibleQuote{我发觉在我的肢体中另有一条法律，与我理智的法律交战，把我掳去，叫我隶属于那在我肢体中罪和死的法律。我真不幸啊！谁能救我脱离这该死的肉身呢？除非是天主的恩宠，藉着我们的主耶稣基督。}(\BibleRef{罗 7:23}) \BibleQuote{世界已经钉在十字架上于我，我于世界。}(\BibleRef{迦 6:14})\par

\Person{奥古斯丁}说：我承认并拥抱神圣圣经的见证，我要用几句话（愿天主赐予恩宠）说明它们如何与我的信仰一致。我说，那最初被造的人有自由意志的运用。他被造得如此完美，以至于如果他愿意遵守天主的诫命，绝对没有任何事物能抗拒他的意志。但是，在他自愿犯罪之后，我们这些从他的后代繁衍下来的人，就被抛入了必然性。但每个人都可以稍加思考发现我所说的是真的。因为在今天的行动中，在我们被任何习惯牵累之前，我们可以自由选择做或不做某事。但是，当我们借着那自由做了某事，而那有害的甜蜜和快乐占据了心灵时，心灵就被自己的习惯如此束缚，以至于此后它无法战胜它因犯罪而为自己塑造的东西。我们看到许多人不想发誓，但因为舌头已经习惯，他们无法阻止那些我们只能归因于邪恶之根的话语从口中说出。为了我可以与你讨论这些话，既然它们没有从你口中消失，但愿它们能被你的心理解：你指着\OriginalTerm{护慰者}{Paraclete}发誓。因此，如果你希望通过经验发现我所说的是否真实，就决心不要发誓。你会看到，那个习惯就沿着它已经习惯的方式前进。这就是那与灵魂作战的东西，在肉身中形成的习惯。这确实是肉身的思念，只要它不能如此服从天主的法律，它就是肉身的思念；但当灵魂被光照，它就不再是肉身的思念了。因为这样说：肉身的思念不能服从天主的法律，就好像说雪不能变热一样。因为只要它是雪，它绝不能变热。但是，正如雪被热量融化而变成温热，同样，肉身的思念——即与肉身形成的习惯——当我们的心智被光照，也就是当天主使整个人服从于神圣法律的抉择时，就代替了灵魂的邪恶习惯，形成了良好的习惯。因此，主关于两棵树——一好一坏——所说的话是最真实的，你已经想起了它们，说它们结出各自的果实；也就是说，好树不能结坏果子，坏树也不能结好果子，但这是当它还是坏的时候。让我们取两个人，一个好人和一个坏人。只要他是好人，他就不能结坏果子；只要他是坏人，他就不能结好果子。但是，为了使你知道这两棵树被主如此放置，以致于自由选择在其中被象征，这两棵树不是本性而是我们的意志，祂自己在福音中说道：\BibleQuote{要么使树好，要么使树坏}\BibleRef{玛 12:35}。谁能创造本性呢？因此，如果我们被命令使树要么好要么坏，那么选择什么是我们的事。因此，关于人的那罪和关于与肉身形成的灵魂的那习惯，宗徒说：\BibleQuote{不要让任何人欺骗你们}\BibleRef{弗 5:6} \BibleQuote{天主所创造的一切都是好的}\BibleRef{弟前 4:4}。你们也引用了的同一位宗徒说：\BibleQuote{正如因一人悖逆，许多人成为罪人；同样，因一人服从，许多人成为义人}\BibleRef{罗 5:19} \BibleQuote{既然死是因一人而来，死人复活也是因一人而来}。因此，只要我们带着属土之人的肖像，即只要我们按肉身生活（这也称为旧人），我们就受我们习惯的必然性束缚，以致不能做我们所愿意的。但是，当天主的恩宠将神圣的爱吹入我们内，并使我们服从祂的旨意时，我们就听到说：\BibleQuote{你们蒙召，是为了自由}\BibleRef{迦 5:13}，以及\BibleQuote{天主的恩宠使我摆脱罪恶与死亡的法律}\BibleRef{罗 8:2}。罪恶的法律是：凡犯罪的，必要死亡。当我们开始成为义人时，我们就从这法律中被解放。死亡的法律是那对人说的：\BibleQuote{你是土，仍要归于土}\BibleRef{创 3:19}。因为我们都是这样出生的，因为我们是土，又因为我们都是这样出生的，因为我们是土，我们都要因第一个人的罪过而归于土。但是，由于天主的恩宠使我们摆脱罪恶与死亡的法律，我们归向正义便得解放；因此，这同一肉身只要我们还留在罪中，就以惩罚折磨我们，而在复活中则服从于我们，并不以任何逆境动摇我们遵守天主的法律和祂的诫命。因此，既然我已回应了你们的问题，请你们俯允按我所愿回答，如果本性是反对天主的，那么罪怎么能归咎于我们？我们并非自愿，而是由那无人能伤害的天主亲自送入那本性中的。\par

\Person{福图纳图斯}说：正如主也对他的门徒们说：\BibleQuote{看，我派遣你们犹如羊进入狼群}\BibleRef{玛 10:16}。因此必须知道，我们的救主并非怀着敌意派遣祂的羔羊，即祂的门徒，进入狼群，除非有一些对立之处，祂要用狼的比喻来指明，祂也把祂的门徒派遣到那里去；使那些可能在狼群中被欺骗的灵魂，能被召回到它们自己的本体。因此也可以看出我们回归的时代和我们岁月的古老，那就是在创世之前，灵魂就这样被派遣去对抗那相反的本性，为使他们通过受难制服它，从而将胜利归还给天主。因为同一位宗徒说，不但要与血肉较量，也要与率领者、掌权者、黑暗世界的霸主以及邪恶的神灵较量。\BibleRef{弗 6:12} 因此，如果两处都居住着邪恶，并被看作邪恶，那么现在邪恶不仅在我们的身体中，而且在整个世界，那里似乎居住着灵魂，它们居住在那穹苍之下并被束缚。\par

23. \Person{奥古斯丁}说：\Quote{主将祂的羔羊送入狼群之中}，即：将义人送入罪人之中，为宣讲在人类时期从不可估量的神圣智慧所接受的福音，好将我们从罪恶中召叫至正义。但宗徒所说，我们的战斗不是对抗血肉，而是对抗率领者与掌权者，以及所引用的其他经文，这表示魔鬼及其天使，正如我们也一样，因罪恶而堕落和失足，并已夺取了属地之物的占有权，即罪人，只要我们仍是罪人，就处于他们的轭下；正如当我们成为义人时，我们将处于正义的轭下；我们与他们战斗，为能转向正义而摆脱他们的统治。因此，也请您屈尊回答我所问的一个问题：天主能否受到伤害？还是不能？但我请求您回答：祂不能受到伤害。\par

\Person{福尔图纳图斯}说：祂不能受到伤害。\par

24. \Person{奥古斯丁}说：那么，按照您的信仰，祂为何派遣我们到此地来呢？\par

\Person{福尔图纳图斯}说：我公开承认的是：天主不能受到伤害，并且祂引导我们到此地来。但既然这与您的观点相悖，请您说说您如何解释灵魂在此地，而我们的天主既愿藉祂的诫命又愿藉祂所派遣的圣子来解放灵魂。\par

25. \Person{奥古斯丁}说：既然我看见您无法回答我的问题，而想问我一些事，看，我让您满意，只要您记住您没有回答我的问题。灵魂为何在此世界处于苦难之中，我不仅刚才，而且之前多次解释过：灵魂犯了罪，因此是痛苦的。它接受了自由意志，按其所愿使用了自由意志；它堕落了，从幸福中被逐出，陷入了苦难。关于这一点，我向您引述了宗徒的见证，他说：\BibleQuote{正如因一人死亡，同样也因一人有了死者的复活。}您还问什么？因此，请您回答：那不能受到伤害的那一位，为何派遣我们到此地来呢？\par

\Person{福尔图纳图斯}说：必须探求原因，为何灵魂来到此地，或为何天主愿意从此地解放生活在邪恶中的灵魂？\par

26. \Person{奥古斯丁}说：我问您这个原因，即：如果天主不能受到伤害，祂为何派遣我们到此地来？\par

\Person{福尔图纳图斯}说：人们问我们，如果邪恶不能伤害天主，为何灵魂被派遣到此地，或出于什么理由它与世界混合？宗徒的话显明了这一点：\BibleQuote{被造物岂能对造它的说：你为什么这样造了我？}\newline{}\BibleRef{罗 9:20}因此，如果必须辩护这个原因，就必须问祂，为何祂派遣灵魂，没有必然性迫使祂。但如果派遣灵魂是出于必然性，那么解放它也有正当的意愿。\par

27. \Person{奥古斯丁}说：那么，天主受必然性所迫吗？\par

\Person{福尔图纳图斯}说：正是如此。请不要对所说的话引致反感，因为我们并不使天主服从必然性，而是使祂自愿地派遣了灵魂。\par

28. \Person{奥古斯丁}说：回想一下上面所说的话。它写道：\Quote{但如果派遣灵魂是出于必然性，那么解放它也有正当的意愿。\Person{奥古斯丁}说：我们听到了：\InnerQuote{但如果派遣灵魂是出于必然性，那么解放它也有正当的意愿。}}因此，您说派遣灵魂是出于必然性。但如果您只想说\Quote{派遣的意愿}，我再补充这一点：那不能受到伤害的那一位，竟有残酷的意愿将灵魂送到如此大的苦难中。因为我说话是为了驳斥这个说法，我恳求祂的仁慈宽恕，我们从祂那里怀有望德，从异端者的一切错误中得到解放。\par

\Person{福尔图纳图斯}说：您断言我们说天主在派遣灵魂时是残酷的，但天主造了人，向他嘘了一口气，赋予他灵魂，这灵魂祂确已预知将卷入未来的苦难，并且由于邪恶而无法恢复其遗产。这要么属于无知者，要么属于将灵魂交付给上述邪恶者。我引用这一点，是因为您不久前说，天主接纳了灵魂，而非灵魂源自祂；因为接纳是另一回事。\par

29. \Person{奥古斯丁}说：关于接纳，我记得几天前根据宗徒的见证说过，宗徒说我们蒙召成为义子的接纳。\BibleRef{弗 1:5}因此，这不是我的回答，而是宗徒的回答。关于这件事，即关于那接纳，如果我们愿意，可以在适当的时候探讨；当您回答了反驳我的论点时，我将毫不迟延地作出回答。\par

\Person{福尔图纳图斯}说：我说灵魂曾出离去对抗一个相反的本性，这个本性不能伤害天主。\par

30. \Person{奥古斯丁}说：那出离有何必要，既然天主不受任何事物伤害，便没有什么要保护的？\par

\Person{福尔图纳图斯}说：您凭良心相信基督来自天主吗？\par

31. \Person{奥古斯丁}说：您又在问我。请回答我的问题。\par

\Person{福尔图纳图斯}说：我在信德中如此领受了，即祂藉天主的旨意来到了此地。\par

32. \Person{奥古斯丁}说：而我说：为什么那全能、不可侵犯、不变、不受任何事物伤害的天主，将灵魂派遣到此地，陷入苦难、错误和我们所遭受的那些事物？\par

\Person{福尔图纳图斯}说：因为经上说过：\BibleQuote{我有权舍掉我的灵魂，也有权再取回它。}\BibleRef{若 10:18}现在祂说灵魂是藉天主的旨意出离的。\par

33. \Person{奥古斯丁}说：我要求理由，为什么天主，既然绝不可能受到任何伤害，却将灵魂派遣到此地？\par

\Person{福尔图纳图斯}说：我们已经说过，天主绝不可能受到任何伤害，并且我们说过灵魂处于一个相反的本性中，因此它给这个相反的本性设定了界限。对相反的本性施加了约束之后，天主便取回这灵魂。因为祂亲自说过：\BibleQuote{我有权舍掉我的灵魂，也有权取回它。}圣父赐给我舍掉灵魂和取回灵魂的能力。因此，那位在天主圣子内说话的天主是指哪一个灵魂呢？显然是指我们的灵魂，它被拘在这些身体中，是出于祂的旨意而来，又出于祂的旨意而被取回。\par

34. \Person{奥古斯定}说：为什么我们的主说：\BibleQuote{我有权舍弃我的生命，也有权再取回它}，这是众人皆知的；因为祂即将受苦并复活。但我一而再、再而三地请问你：如果天主不可能以任何方式遭受损害，那祂为何还要将灵魂遣送到这里来？\par

\Person{福尔图纳图斯}说：为了给\OriginalTerm{反性}{contrary nature}设定一个界限。\par

35. \Person{奥古斯定}说：难道全能的、仁慈的、至高的天主，为了约束\OriginalTerm{反性}{contrary nature}，竟然愿意让它受到限制，好使我们不受约束吗？\par

\Person{福尔图纳图斯}说：但祂正是在召唤我们回归祂自己。\par

36. \Person{奥古斯定}说：如果祂是从一种不受约束的状态、从罪恶、从错误、从痛苦中召唤我们回归，那么，灵魂为何必须经历如此巨大的苦难，历经如此漫长的时期，直到世界的终结？既然您说那位遣送灵魂的天主不可能以任何方式遭受损害。\par

\Person{福尔图纳图斯}说：那我该说什么呢？\par

37. \Person{奥古斯定}说：我知道您无话可说，而且我本人，当我在你们中间时，也从未在这个问题上找到任何可说的话，正是因此，我从上受到警醒，得以离开那个错误，并归向\OriginalTerm{天主教信仰}{Catholic faith}，更确切地说，是通过那位的宽仁重新记起它，祂不允许我永远固守在这谬误之中。但如果您承认自己无可答复，那么，我将向所有聆听并探求的人阐述天主教信仰，既然他们是信者，只要他们允许并愿意。\par

\Person{福尔图纳图斯}说：在不损害我所宣认的前提下，我可以说：当我与我的上司们重新审议您所提出的反对意见后，如果他们无法回应我的这个问题——即您现在同样向我提出的——那么，我将考虑（因为我渴望我的灵魂能通过一个确定的信仰获得释放）前来探讨您向我提议并承诺将要展示的这件事。\par

\Person{奥古斯定}说：感谢天主。\par

\SourceRecord{Translated by Albert H. Newman.；Nicene and Post-Nicene Fathers, First Series；Vol. 4.；Edited by Philip Schaff.；Buffalo, NY: Christian Literature Publishing Co.,；1887.；<http://www.newadvent.org/fathers/1404.htm>.}
